\documentclass[twoside,english,brazilian]{UNISINOSmonografia}
\usepackage[utf8]{inputenc} % charset do texto (utf8, latin1, etc.)
\usepackage[T1]{fontenc}    % encoding da fonte (afeta a sep. de sílabas)
\usepackage[alf]{abntex2cite}
\usepackage{booktabs}       % para tabelas mais elegantes
\usepackage{array}          % para formatação avançada de tabelas
\usepackage{longtable}      % para tabelas que quebram páginas
\usepackage[alf]{abntex2cite} % autor-data ABNT
\usepackage{graphicx}
\usepackage{caption}
\usepackage{subcaption} % opcional, para subfiguras

%=======================================================================
% Dados gerais sobre o trabalho (mantidos para capa/folha de rosto).
%=======================================================================
\autor{Figueiredo Carolino}{Gabriel}
\titulo{Sistema Inteligente de Monitoramento e Irrigação Automatizada para Plantas Domésticas baseado em IoT}
\subtitulo{Solução Integrada de Automatização Residencial com ESP32, Firebase e Interface Web Adaptativa}
\orientador[Dr.]{Lacerda}{Guilherme}

\unidade{Unidade Acadêmica de Graduação}
\curso{Curso de Bacharelado em Ciência da Computação}
\natureza{%
Monografia apresentada como requisito parcial para obtenção do título de Bacharel em Ciência da Computação, pelo Curso de Ciência da Computação da Universidade do Vale do Rio dos Sinos (UNISINOS)
}
\local{Porto Alegre}
\ano{2025}

% Palavras-chave em PT para metadados (atualizadas)
\palavrachave{brazilian}{automação residencial}
\palavrachave{brazilian}{Internet das Coisas}
\palavrachave{brazilian}{monitoramento inteligente}
\palavrachave{brazilian}{ESP32}
\palavrachave{brazilian}{Firebase}
\palavrachave{brazilian}{plantas domésticas}

%=======================================================================
% Início do documento.
%=======================================================================
\begin{document}
\capa
\folhaderosto

%=======================================================================
% Página inicial do artigo no padrão do exemplo da UNISINOS
% (título PT/EN, autoria com notas de rodapé, resumo/abstract e palavras-chave).
%=======================================================================
\clearpage
\thispagestyle{plain}
\begin{center}
\MakeUppercase{TÍTULO EM PORTUGUÊS: Sistema Inteligente de Monitoramento e Irrigação Automatizada para Plantas Domésticas baseado em IoT}\\[0.75em]
\MakeUppercase{TÍTULO EM OUTRO IDIOMA: Intelligent Monitoring and Automated Irrigation System for Household Plants based on IoT}\\[1.5em]

Gabriel Figueiredo Carolino\footnote{Bacharelando em Ciência da Computação (UNISINOS). E-mail: \texttt{gabrielcarolino962@gmail.com}.} \\
Orientador: Prof. Dr. Guilherme Lacerda\footnote{UNISINOS. E-mail: \texttt{guilhermeslacerda@gmail.com}.}
\end{center}

\bigskip

\noindent\textbf{Resumo:} Este artigo apresenta o desenvolvimento e avaliação de um sistema inteligente de monitoramento e irrigação automatizada para ambientes residenciais, focado na facilidade de uso e conveniência para cuidadores de plantas domésticas. A solução integra tecnologias de Internet das Coisas (IoT) através de sensores ambientais múltiplos conectados a um microcontrolador ESP32, comunicação em tempo real via Firebase Realtime Database e uma interface web responsiva com dois modos de operação. O sistema elimina a necessidade de monitoramento manual constante, prevenindo tanto a morte de plantas por falta de água quanto problemas causados por irrigação excessiva. A arquitetura proposta adapta-se automaticamente às necessidades específicas de diferentes espécies vegetais e tamanhos de vasos, proporcionando tranquilidade e praticidade aos usuários através de automação residencial inteligente.

\medskip
\noindent\textbf{Palavras-chave:} automação residencial; Internet das Coisas; monitoramento inteligente; ESP32; Firebase; plantas domésticas.

\medskip
\noindent\textbf{Abstract:} This paper presents the development and evaluation of an intelligent monitoring and automated irrigation system for residential environments, focused on ease of use and convenience for household plant caregivers. The solution integrates Internet of Things (IoT) technologies through multiple environmental sensors connected to an ESP32 microcontroller, real-time communication via Firebase Realtime Database, and a responsive web interface with two operation modes. The system eliminates the need for constant manual monitoring, preventing both plant death from lack of water and problems caused by excessive irrigation. The proposed architecture automatically adapts to the specific needs of different plant species and pot sizes, providing peace of mind and practicality to users through intelligent home automation.

\medskip
\noindent\textbf{Key-words:} home automation; Internet of Things; intelligent monitoring; ESP32; Firebase; household plants.

%=======================================================================
% 1 INTRODUÇÃO
%=======================================================================
\chapter{Introdução}

A Internet das Coisas (IoT) tem experimentado um crescimento exponencial nos últimos anos, com previsões indicando que o número de dispositivos conectados alcançará 25 bilhões até 2025, representando um mercado global estimado em US\$ 1,1 trilhão \cite{pacete2022iot}. No segmento específico de automação residencial, o mercado de smart homes cresceu 32\% em 2022, atingindo um valor de US\$ 80 bilhões globalmente, com projeções de chegar a US\$ 338 bilhões até 2030 \cite{fbi2025_smart_home_sa}. Este crescimento reflete uma demanda crescente por soluções tecnológicas que simplifiquem tarefas domésticas e melhorem a qualidade de vida dos usuários.

% IMAGEM 1: Gráfico de crescimento do mercado IoT e automação residencial
% DESCRIÇÃO: Gráfico de linhas mostrando dados estatísticos 2020-2030 com duas curvas:
% - Mercado IoT global (bilhões de dispositivos conectados)
% - Mercado de automação residencial (bilhões de dólares)
% Incluir segmentação por aplicação (irrigação, climatização, segurança)
% Destacar a projeção de crescimento exponencial e oportunidades de mercado
% Fonte: Dados compilados de pacete2022iot e fbi2025_smart_home_sa
\begin{figure}[htbp]
\centering
% \includegraphics[width=0.8\textwidth]{figuras/crescimento_mercado_iot.png}
\caption{Projeção de crescimento do mercado IoT e automação residencial (2020-2030)}
\label{fig:crescimento_mercado}
\end{figure}

O cultivo de plantas em ambientes domésticos tem se tornado cada vez mais popular, especialmente em contextos urbanos onde a conexão com a natureza é limitada. No entanto, o cuidado adequado de plantas domésticas apresenta desafios significativos para a maioria das pessoas. A irrigação adequada representa o principal obstáculo, exigindo conhecimento específico sobre as necessidades hídricas de cada espécie, monitoramento constante das condições do solo e disponibilidade regular para realizar os cuidados necessários.

A rotina moderna, com suas demandas profissionais e pessoais intensas, frequentemente resulta em esquecimento ou negligência no cuidado das plantas. Muitas pessoas relatam frustração ao perder plantas por falta de água durante viagens, períodos de trabalho intenso ou simplesmente por não conseguir estabelecer uma rotina consistente de cuidados. Por outro lado, o excesso de zelo pode levar à irrigação excessiva, causando apodrecimento das raízes e morte das plantas por motivos opostos.

A automatização residencial, potencializada pelas tecnologias de Internet das Coisas (IoT), oferece uma solução elegante para esses desafios. Sistemas inteligentes podem assumir a responsabilidade pelo monitoramento contínuo e cuidados automatizados, liberando os usuários da preocupação constante e garantindo que as plantas recebam exatamente a quantidade de água necessária no momento adequado.

Este trabalho aborda especificamente a necessidade de sistemas de automatização que sejam ao mesmo tempo tecnicamente robustos e extremamente simples de usar. A solução proposta reconhece que nem todos os usuários possuem conhecimento técnico avançado, oferecendo interfaces adaptativas que atendem tanto iniciantes quanto usuários experientes. O foco central está na tranquilidade e praticidade proporcionadas pela automação, permitindo que as pessoas desfrutem dos benefícios das plantas domésticas sem as preocupações tradicionais associadas aos seus cuidados.

\medskip
\noindent\textbf{Problema Central.} Como desenvolver um sistema de automação residencial que elimine efetivamente as preocupações e dificuldades relacionadas ao cuidado de plantas domésticas, proporcionando monitoramento inteligente e intervenções automatizadas que se adaptem às necessidades específicas de diferentes espécies e condições ambientais?

\medskip
\noindent\textbf{Objetivo Geral.} Desenvolver e validar um sistema IoT integrado que automatize completamente o monitoramento e irrigação de plantas domésticas, oferecendo uma solução prática e confiável que permita aos usuários cultivar plantas com sucesso independentemente de seu nível de experiência ou disponibilidade de tempo.

\medskip
\noindent\textbf{Objetivos Específicos.}
\begin{itemize}
    \item Implementar monitoramento ambiental contínuo através de sensores de umidade do solo, temperatura, umidade do ar e luminosidade integrados a microcontrolador ESP32;
    \item Desenvolver sistema de controle automatizado com algoritmos adaptativos que respondam às necessidades específicas de diferentes espécies de plantas;
    \item Criar interface web responsiva com dois modos de operação (simples e avançado) para atender usuários com diferentes níveis de conhecimento técnico;
    \item Estabelecer comunicação em tempo real através de Firebase Realtime Database para monitoramento remoto e controle via dispositivos móveis e computadores;
    \item Validar a eficácia do sistema através de testes práticos com três vasos de diferentes tamanhos, demonstrando adaptabilidade e confiabilidade;
    \item Desenvolver metodologia de calibração específica para diferentes volumes de substrato, considerando as particularidades de vasos pequenos, médios e grandes.
\end{itemize}

\noindent\textbf{Contribuições Principais.} Este trabalho oferece as seguintes contribuições para a área de automação residencial e IoT aplicada ao cultivo doméstico:

A primeira contribuição consiste no desenvolvimento de uma arquitetura IoT completa e acessível que integra hardware de baixo custo (ESP32 e sensores) com serviços em nuvem (Firebase) e interface web moderna. Esta arquitetura demonstra como tecnologias emergentes podem ser aplicadas de forma prática para resolver problemas cotidianos reais.

A segunda contribuição refere-se à criação de interfaces adaptativas que democratizam o acesso à tecnologia de automação. O sistema oferece um modo simples com controles intuitivos para usuários iniciantes, enquanto mantém funcionalidades avançadas disponíveis para usuários experientes, eliminando barreiras técnicas sem sacrificar capacidades.

A terceira contribuição apresenta uma metodologia científica para calibração de sensores capacitivos de umidade considerando diferentes volumes de substrato. Esta metodologia é especialmente relevante pois reconhece que vasos de tamanhos diferentes requerem calibrações específicas devido às variações na distribuição de umidade e características físicas do substrato.

Além de uma implementação prática, o trabalho contribui metodologicamente ao propor uma forma sistematizada de calibração e validação para sistemas IoT domésticos, estabelecendo diretrizes que podem ser aplicadas em contextos similares. 

Finalmente, o trabalho contribui com uma análise detalhada dos fatores que influenciam a adoção de tecnologias de automação residencial, identificando barreiras práticas e propondo soluções que facilitem a implementação em larga escala.

\noindent\textbf{Organização do Trabalho.} Este trabalho está organizado da seguinte forma: na Seção 2, é apresentada a fundamentação teórica necessária, abordando conceitos de automação residencial, Internet das Coisas e tecnologias habilitadoras, seguida de análise detalhada de trabalhos relacionados na área. A Seção 3 descreve a metodologia científica adotada, caracterizada como pesquisa de design science, e detalha os materiais utilizados e procedimentos experimentais. A Seção 4 apresenta a arquitetura completa do sistema proposto, incluindo aspectos de hardware, software e integração entre componentes. A Seção 5 documenta a implementação prática e os resultados obtidos nos testes de validação, incluindo análise de desempenho técnico e experiência do usuário. A Seção 6 discute as implicações dos resultados, limitações identificadas e propõe direções para trabalhos futuros. Finalmente, a Seção 7 apresenta as conclusões e reflexões sobre o impacto e contribuições da solução desenvolvida.

%=======================================================================
% 2 FUNDAMENTAÇÃO TEÓRICA E TRABALHOS RELACIONADOS
%=======================================================================
\chapter{Fundamentação Teórica e Trabalhos Relacionados}

\section{Automação Residencial e Qualidade de Vida}

A automação residencial representa uma evolução natural na busca por maior conforto, segurança e eficiência nos ambientes domésticos. Diferentemente dos primeiros sistemas automatizados que focavam principalmente em controle de temperatura e iluminação, as soluções modernas abraçam uma visão holística do lar como um ecossistema inteligente e responsivo às necessidades dos habitantes.

No contexto específico do cultivo doméstico, a automação assume um papel particularmente importante ao eliminar uma fonte significativa de ansiedade e responsabilidade para muitas pessoas. A preocupação constante com o bem-estar das plantas, especialmente durante ausências prolongadas, representa um fator limitante que impede muitos indivíduos de desfrutar dos benefícios reconhecidos das plantas domésticas, incluindo melhoria da qualidade do ar, redução do estresse e conexão com a natureza.

Estudos na área de psicologia ambiental demonstram que a presença de plantas em ambientes internos contribui significativamente para o bem-estar psicológico e físico dos ocupantes \cite{fbi2025_smart_home_sa}. No entanto, essas mesmas pesquisas identificam que a "ansiedade de cuidado" – a preocupação constante sobre se as plantas estão recebendo cuidados adequados – pode paradoxalmente reduzir os benefícios psicológicos esperados. A automação resolve essa contradição ao assumir a responsabilidade pelo cuidado técnico, permitindo que as pessoas se concentrem nos aspectos prazerosos da convivência com plantas.

A Internet das Coisas (IoT) fundamenta-se na premissa de que objetos cotidianos podem se tornar "inteligentes" através da incorporação de sensores, processamento local e conectividade. No ambiente doméstico, essa inteligência distribuída cria oportunidades para sistemas que aprendem, adaptam-se e respondem automaticamente às condições em constante mudança.

Para aplicações de cultivo doméstico, a IoT oferece vantagens específicas que transcendem a simples automação. A capacidade de coleta contínua de dados ambientais permite identificar padrões sutis que seriam imperceptíveis ao monitoramento humano ocasional. Por exemplo, variações graduais na umidade do solo ao longo de semanas podem indicar mudanças nas necessidades da planta devido ao crescimento, alterações sazonais ou desenvolvimento do sistema radicular.

A conectividade proporcionada pela IoT também facilita o monitoramento remoto, aspecto crucial para pessoas que viajam frequentemente ou trabalham longas jornadas. A tranquilidade de poder verificar o status das plantas e, se necessário, intervir remotamente, remove uma barreira significativa para a adoção de plantas domésticas por pessoas com estilos de vida mais dinâmicos.

Além disso, sistemas IoT bem projetados podem acumular dados históricos que se tornam progressivamente mais valiosos, permitindo otimizações baseadas no comportamento específico de cada planta e ambiente. Esta capacidade de "aprendizado" através de dados históricos diferencia soluções IoT de sistemas de automação tradicionais mais simples.

\section{Tecnologias Habilitadoras}

\subsection{Microcontrolador ESP32 e Capacidades de Processamento Local}

O ESP32 representa uma evolução significativa em microcontroladores para aplicações IoT, combinando processamento dual-core, conectividade Wi-Fi e Bluetooth integradas, e consumo energético otimizado em um único chip de baixo custo. Para aplicações de automação residencial, essas características traduzem-se em capacidades práticas importantes.

O processamento dual-core permite que o ESP32 execute simultaneamente tarefas de coleta de dados dos sensores e comunicação com serviços na nuvem, sem comprometer a responsividade do sistema. 

A conectividade Wi-Fi integrada elimina a necessidade de componentes adicionais e simplifica significativamente a integração com redes domésticas existentes. A capacidade de reconexão automática e gerenciamento transparente de credenciais de rede contribui para a confiabilidade operacional necessária em sistemas de automação residencial.

O baixo consumo energético do ESP32, especialmente quando utiliza modos de economia de energia, torna viável a operação contínua por longos períodos. 

\subsection{Sensores Capacitivos e Precisão de Medição}

Sensores capacitivos de umidade do solo representam um avanço significativo em relação aos sensores resistivos tradicionais, oferecendo maior durabilidade, precisão e estabilidade ao longo do tempo. O princípio de funcionamento baseado na medição de constante dielétrica do substrato elimina problemas de corrosão eletrolítica que afetam sensores resistivos.

Para aplicações de cultivo doméstico, a precisão e confiabilidade dos sensores capacitivos são especialmente importantes devido à natureza crítica das decisões de irrigação. Leituras imprecisas podem resultar em irrigação inadequada, comprometendo a saúde das plantas e contradizendo o objetivo fundamental do sistema de automação.

\subsection{Firebase Realtime Database e Sincronização em Tempo Real}

O Firebase Realtime Database oferece capacidades de sincronização de dados em tempo real que são fundamentais para a experiência de usuário em aplicações de monitoramento remoto. A arquitetura push-based garante que atualizações sejam propagadas imediatamente para todas as interfaces conectadas, eliminando a necessidade de atualizações manuais ou polling periódico.

Para sistemas de automação residencial, essa sincronização instantânea é crucial para manter os usuários informados sobre mudanças no status das plantas, especialmente durante eventos críticos como início ou fim de irrigação. A capacidade de visualizar mudanças em tempo real aumenta a confiança no sistema e proporciona transparência operacional.

\section{Calibração de Sensores e Variações de Substrato}

Um aspecto frequentemente subestimado em sistemas de irrigação automatizada é a influência significativa do volume e tipo de substrato nas leituras dos sensores de umidade. Esta questão assume particular importância em aplicações domésticas, onde vasos de diferentes tamanhos e tipos de terra são comuns.

A calibração de sensores capacitivos deve considerar que a distribuição de umidade em substratos varia significativamente com o volume do recipiente. Em vasos pequenos, a umidade distribui-se mais uniformemente e responde rapidamente à irrigação. Em vasos grandes, gradientes de umidade podem existir entre diferentes regiões, e a resposta à irrigação é mais lenta e heterogênea.

Além disso, diferentes tipos de substrato apresentam propriedades dielétricas distintas, afetando diretamente as leituras de sensores capacitivos. Substratos com maior componente orgânico, diferentes granulometrias ou aditivos como vermiculita ou perlita podem requerer ajustes específicos na calibração.

Esta realidade prática exige que sistemas de automação residencial incorporem flexibilidade na calibração, permitindo ajustes específicos para cada combinação de vaso, substrato e espécie de planta. A metodologia de calibração deve ser suficientemente simples para usuários não técnicos, mas precisa o suficiente para garantir operação confiável.

\section{Trabalhos Relacionados e Análise Comparativa}

A literatura científica apresenta diversos trabalhos que abordam automação de irrigação em diferentes contextos e escalas. \citeonline{nunes_2025_jardim_inteligente} desenvolveram especificamente um sistema para plantas domésticas, com ênfase na simplicidade de uso. No entanto, o trabalho não aborda adequadamente questões de calibração para diferentes tipos de vasos e substratos, limitando sua aplicabilidade prática.

Uma revisão abrangente conduzida por \citeonline{lima2019} analisa mais de 40 trabalhos na área de automação de irrigação, identificando tendências e lacunas. Os autores destacam que a maioria dos trabalhos foca em aspectos técnicos de hardware e comunicação, com atenção limitada às necessidades práticas dos usuários finais e questões de usabilidade.

Outro trabalho relevante é o de \citeonline{volkmann2022}, que apresenta um sistema de irrigação automatizada de baixo custo utilizando MQTT e interface móvel. Apesar de avançar na integração com dispositivos móveis, o estudo concentra-se em ambientes de jardins e não aborda a adaptação para diferentes volumes de vasos domésticos, além de não considerar calibração específica para substratos variados.

\subsection{Análise Comparativa Detalhada}

Para proporcionar uma visão mais clara do posicionamento deste trabalho em relação ao estado da arte, a Tabela \ref{tab:comparativo} apresenta uma análise comparativa dos principais trabalhos relacionados, destacando tecnologias utilizadas, limitações identificadas e como a presente proposta se diferencia.

\begin{table}[htbp]
\centering
\caption{Análise comparativa de trabalhos relacionados em automação de irrigação}
\label{tab:comparativo}
\scriptsize
\begin{tabular}{p{2.2cm}p{2.0cm}p{2.6cm}p{2.8cm}p{3.0cm}}
\toprule
\textbf{Autor/Ano} & \textbf{Tecnologias} & \textbf{Aplicação} & \textbf{Limitações} & \textbf{Diferencial da Nossa Proposta} \\
\midrule
Nunes (2025) & ESP8266, sensor capacitivo & Plantas domésticas & Calibração única; ignora tamanhos de vaso & Calibração específica por volume; 3 tamanhos testados \\
\midrule
Lima (2019) & Revisão: Arduino, ESP & Agricultura familiar & Foco técnico; pouca atenção ao usuário & Centrado no usuário; validação UX com TAM \\
\midrule
Volkmann (2022) & MQTT, interface móvel & Agricultura de precisão & Sem calibração por vaso; foco em campo & Adaptação doméstica; calibração por substrato/volume \\
\midrule
\textbf{Nossa Proposta} & \textbf{ESP32, Firebase, multi-sensores} & \textbf{Plantas domésticas} & \textbf{1 vaso por dispositivo} & \textbf{Baixo custo + usabilidade + validação científica} \\
\bottomrule
\end{tabular}
\end{table}

A análise comparativa evidencia que, embora existam diversos trabalhos na área, cada um apresenta limitações específicas que restringem sua aplicabilidade prática no contexto doméstico. O presente trabalho busca preencher essas lacunas através de uma abordagem integrada que equilibra sofisticação técnica com simplicidade de uso, além de validar cientificamente a solução em cenários realísticos.

\subsection{Lacunas Identificadas na Literatura}

A análise sistemática da literatura revela várias lacunas importantes que este trabalho busca abordar, proporcionando contribuições significativas para o avanço da área:

\textbf{Primeira lacuna}: Existe uma carência notável de sistemas que equilibrem efetivamente sofisticação técnica com simplicidade de uso. A maioria dos trabalhos acadêmicos presume usuários com conhecimento técnico significativo, criando barreiras de adoção para o público geral. Esta lacuna é particularmente problemática considerando que o objetivo fundamental da automação residencial é simplificar a vida das pessoas, não adicionar complexidade.

\textbf{Segunda lacuna}: A questão da calibração específica por tipo e tamanho de vaso é frequentemente negligenciada ou tratada superficialmente na literatura existente. Esta limitação compromete significativamente a precisão operacional em aplicações domésticas reais, onde a diversidade de vasos e substratos é a norma, não a exceção. A falta de metodologias sistematizadas para calibração adaptativa representa uma barreira técnica importante para a implementação prática desses sistemas.

\textbf{Terceira lacuna}: Poucos trabalhos abordam adequadamente aspectos psicológicos e comportamentais da automação residencial, como a importância da transparência operacional e controle do usuário para aceitação da tecnologia. A literatura técnica tende a focar em métricas de desempenho hardware e software, negligenciando fatores humanos que são críticos para o sucesso de sistemas destinados ao uso doméstico.

\textbf{Quarta lacuna}: Existe uma oportunidade significativa na criação de interfaces verdadeiramente adaptativas que atendam tanto usuários iniciantes quanto experientes sem sacrificar funcionalidade ou aumentar complexidade desnecessariamente. A maioria das soluções existentes opta por interfaces únicas que tendem a ser ou muito simples (limitando funcionalidade) ou muito complexas (intimidando usuários iniciantes).

\textbf{Quinta lacuna}: A validação experimental na literatura existente raramente considera cenários realísticos de uso doméstico, com muitos trabalhos limitando-se a testes de laboratório ou validações técnicas isoladas. 

Este trabalho contribui para preencher essas lacunas através de uma abordagem metodologicamente rigorosa que combina robustez técnica com foco centrado no usuário, atenção específica a aspectos práticos de implementação doméstica, e validação experimental abrangente que considera tanto aspectos técnicos quanto experiência do usuário em condições realísticas de uso.

O diferencial metodológico desta pesquisa reside na integração sistemática de múltiplas perspectivas – técnica, de usabilidade e de validação experimental – em uma solução coesa que demonstra como tecnologias IoT podem ser efetivamente aplicadas para resolver problemas cotidianos reais, mantendo a simplicidade como princípio orientador sem comprometer capacidades técnicas essenciais.

\end{document}
